\documentclass{article}
\usepackage[margin=0.8in]{geometry}
\usepackage{amsmath, amssymb, amsthm}
\usepackage{hyperref}
\usepackage{enumerate}
\usepackage{array}

\newtheorem{proposition}{Proposition}
\newtheorem{definition}{Definition}
\newtheorem{theorem}{Theorem}
\newtheorem{remark}{Remark}

\title{Introduction to University Mathematics}
\author{}
\date{}

\begin{document}
\maketitle
\tableofcontents
\newpage

\section{Lecture 1}
\subsection{The natural numbers and induction}

\begin{definition}
  \normalfont
  A \underline{natural number} is a member of the sequence 0, 1, 2, ... formed by starting from 0 and successively adding 1. Set of all natural number is written as $\mathbb{N} = \{1,2,...\}$.
\end{definition}

\begin{remark}
  Inclusion of 0 in natural number is not universal. Sometimes it is not included.
\end{remark}

Natural number can be added and multiplied to produce natural numbers . i.e. If $m,n\in\mathbb{N}$ then $m+n,m*n\in\mathbb{N}$ ($m*n$ is often written as $mn$). 

0 and 1 are two special natural numbers because they are additive and multiplicative \underline{identities}. i.e. $\forall n\in\mathbb{N}, n + 0 = n \text{ and } n * 1 = n$. The natural numbers have ordering so it makes sense to have the notion comparison between two natural numbers. All the mathematical objects do not enjoy this property (for instance matrices).

\begin{definition}
  \normalfont
  Let $m, n \in\mathbb{N}$. We write $m\leq n$ to mean that $\exists k \in \mathbb{N}:m+k = n$.
\end{definition}

\begin{theorem}
  \normalfont
  (Principle of mathematical induction) Let $P(n)$ be a family of statements indexed by the natural numbers. Suppose:
  \begin{enumerate}[i.]
    \item $P(0)$ is true, and
    \item $\forall n\in\mathbb{N}$, if $P(n)$ and $P(n+1)$ is true
  \end{enumerate}
  then $P(n)$ is true for all $n\in\mathbb{N}$.
\end{theorem}

\begin{proposition}
  \normalfont
  $\forall n \in N$, \[\sum_{k=0}^{n}k = \frac{n(n+1)}{2}.\]
\end{proposition}

\begin{proof}
  $P(0)$ is true since for $n = 0$, LHS = RSH = 0.

  Suppose $P(n)$ is true, then
  $$
  \sum^{n+1}_{k=0} = \sum^{n}_{k=0}k + (n + 1) = \frac{n(n+1)}{2} + (n+1) = \frac{(n+1)(n+2)}{2}
  $$

  So $P(n+1)$ is also true. Therefore by induction, $P(n)$ is true for all $n\in N$
\end{proof}

\begin{theorem}
  \normalfont
  (Strong induction) Let $P(n)$ be a family of statements indexed by $N$. Suppose
  \begin{enumerate}[i.]
    \item $P(0)$ is true, and 
    \item $\forall n\in \mathbb{N}$, if $P(0), P(1),\, ...\,, P(n)$ are true,
  \end{enumerate}
  then $P(n+1)$ i true. This implies $P(n)$ is true $\forall n\in N$.
\end{theorem}

\begin{proof}
  Define $Q(n)$ to be te statement `$P(k)$ is true for $k=0,1,,\, ..., n$'. Then we know that
  \begin{enumerate}[i.]
    \item $Q(0)$ is true, and
    \item $\forall n\in \mathbb{N}$, if $Q(n)$ is true, then $Q(n+1)$ is true.
  \end{enumerate}
  So, by induction, $Q(n)$ is true for all $n\in \mathbb{N}$. Hence $P(n)$ is true for all $n\in\mathbb{N}$.
\end{proof}

\begin{proposition}
  Every natural number greater than 1 can be expressed as a product of one or more primes.
\end{proposition}

\begin{proof}

  Let $P(n)$ be the statement that $n$ can be expressed ad a product of primes. $P(2)$ is true since 2 is a prime itself.

  Let $n>2$, and suppose that $P(n)$ holds for all $m<n$. If $n$ is prime, then $P(n)$ is true. If $n$ is not prime, then $n = rs$ for some $r,s \in\mathbb{N}$, with $r,s < n$. By inductive hypothesis $r$ and $s$ can be expressed as products of primes, and hence so can $n = rs$. So $P(n)$ is true.

  By strong induction, P(n) holds for all $n\in\mathbb{N}$.
\end{proof}

\subsection{Addition and well-ordering principle of $\mathbb{N}$}

\begin{definition}
  (Addition of $\mathbb{N}$) Addition ($+$) is an operation that $\forall m\in \mathbb{N}$,
  
  \begin{enumerate}[i.]
    \item $m + 0 = m$
    \item $\forall n\in \mathbb{N}, m + (n + 1) = (m + n) + 1$
  \end{enumerate}
\end{definition}

\begin{proposition}
  Addition is \underline{associative}, i.e. $\forall x, y, z \in\mathbb{N}$
  
  \begin{equation}
    \label{eqn_associative}
    x + (y + z) = (x + y) + z
  \end{equation}
\end{proposition}

\begin{proof}
  For $z = 0$,

  \begin{equation*}
    \begin{split}
      \textnormal{LHS} &= x + (y + 0)\quad\quad[using (i) from definition]\\
      &= (x+y) + 0\\
      &= \textnormal{RHS}\\
    \end{split}
  \end{equation*}

  Suppose Eqn\,.(\ref{eqn_associative}) holds for $z = n$. Then for $z = n+1$,

  \begin{equation*}
    \begin{split}
      \textnormal{LHS} &= x + (y + (n + 1))\quad\quad[using (ii) from definition]\\
      &= (x + (y + n)) + 1\\
      &= ((x + y) + n) + 1\\
      &= (x + y) + (n + 1)\\
      &= \textnormal{RHS}\\
    \end{split}
  \end{equation*}

  We see that Eqn\,.(\ref{eqn_associative}) holds for $z=n+1$. Therefore, by induction, it holds for all $z\in\mathbb{N}$.
\end{proof}

\begin{proposition}
  (Well-ordering property of the natural numbers)
  Evert non-empty subset of $\mathbb{N}$ has a least element.
\end{proposition}

\begin{proof}
  Assume, for a contradiction, that $S$ is a non-empty subset of $\mathbb{N}$ that does not have a least element. Consider $S^{*} = \{n\in\mathbb{N} : n \notin S\}$.

  Note that $0\in S^{*}$, since  $0\notin S^{*}$ else it would be the least element by default.

  if $0,1,\,...\,,n\in S^{*}$, then $n+1\notin S$, else it could be the lease element. So $n+2\in S^{*}$. By strong induction, $n\in S^{*}$ for all $n\in N$. Hence $S$ is empty, a contradiction.
\end{proof}

\section{Lecture 2}
\subsection{The binomial theorem}

\begin{definition}
  For $n,k\in \mathbb{N}$, we define \underline{binomial coefficient} as

  $$
  \binom{n}{k} = {}^{n}C_{k} = \frac{n!}{(n-k)!k!}\quad\textnormal{for } 0\leq k\leq n\quad\left[\binom{n}{k}=0\textnormal{ for }k > n\right]
  $$

  Note: $0! = 1$ 
\end{definition}

These appear in Pascal's triangle

\begin{center}
  \begin{tabular}{>{$n=}l<{$\hspace{12pt}}*{13}{c}}
    0 &&&&&&&1&&&&&&\\
    1 &&&&&&1&&1&&&&&\\
    2 &&&&&1&&2&&1&&&&\\
    3 &&&&1&&3&&3&&1&&&\\
  \end{tabular}
\end{center}

% lemma
\begin{theorem}
Let $n,k\in\mathbb{N}$ with $1\geq k\geq n$, then
$\binom{n}{k-1} + \binom{n}{k} = \binom{n+1}{k}$
\end{theorem}

\begin{proof}
  \begin{equation*}
    \begin{split}
      \textnormal{LHS} &= \frac{n!}{(n-k+1)!\,(k-1)!} + \frac{n!}{(n-k)!\,k!}\\
      &= \frac{n!\,(k+n-k+1)}{(n-k+1)!\,k!}\\
      &= \frac{(n+1)!}{(n+1-k)!\,k!}\\
      &= \textnormal{RHS}\\
    \end{split} 
  \end{equation*}
\end{proof}

\begin{theorem}
  (Binomial theorem) Let $x, y$ be real or complex numbers and let $n\in\mathbb{N}$. Then
  $$
    (x+y)^{n} = \sum_{k=0}^{n}\binom{n}{k}x^{k}y^{n-k}
  $$
\end{theorem}

\begin{proof}
  For $n=0$, LHS = 1 \& RHS = 1, so true for $n=0$. Suppose it is true for $n$, and consider $n+1$,
\begin{equation*}
  \begin{split}
    (x+y)^{n+1} &=(x+y)(x+y)^{n}\\
    &= (x+y)\sum^{n}_{k=0}\binom{n}{k}x^{k}y^{n-k}\quad\quad\textnormal{[inductive hypothesis]}\\
    &= \sum^{n}_{k=0}\binom{n}{k}x^{k+1}y^{n-k} + \sum^{n}_{k=0}\binom{n}{k}x^{k}y^{n+1-k}\\
    &= x^{n+1} + \sum^{n-1}_{\hat{k}=0,\,\hat{k}=k-1}\binom{n}{\hat{k}}x^{\hat{k}+1}y^{n-\hat{k}} + \sum^{n}_{k=1}\binom{n}{k}x^{k}y^{n+1-k} + y^{n+1}\\
    &= x^{n+1} + \sum^{n}_{k=1}\left[\binom{n}{n-1}x^{k}y^{n+1-k} + \binom{n}{k}x^{k}y^{n+1-k}\right] + y^{n+1}\\
    &= x^{n+1} + \sum^{n}_{k=1}\binom{n+1}{n}x^{k}y^{n+1-k} + y^{n+1}\quad\quad\textnormal{[using the lemma]}\\
    &= \sum_{k=0}^{n+1}\binom{n+1}{k}x^{k}y^{n+1-k}\\
    &= \textnormal{RHS}\\
  \end{split}
\end{equation*}
\end{proof}

\subsection{Introduction to set theory}
A set is a collection of objects. The objects are called the \underline{elements} or \underline{members}. We write the set with elements $a_{1}, a_{2},\,...\,,a_{n}$ as ${a_{1}, a_{2},\,...\,,a_{n}}$.

\begin{definition}
  A is a \underline{subset} of S $(A\subset S)$ if every element of A is an element of S. If $A\neq S$, A is called a \underline{proper subeset} of $S$ $(A\subseteq S)$.
\end{definition}

\begin{definition}
  The empty set $\empty$ is the set with no elements.
\end{definition}

\underline{Examples}:
\begin{itemize}
  \item $\{n\in\mathbb{N}:n \textnormal{divisible by 2}\}\subseteq\mathbb{N}$ is the set of even natural numbers.
  \item $\mathbb{Z}$ is the set of integers $\{0,\pm 1\pm 2,\,...\}$.
  \item $\mathbb{Q}$ is the set of natural numbers $\{\frac{m}{n}: m,n\in\mathbb{Z}, n>0\}$.
  \item $\mathbb{R}$ is the set of real numbers.
  \item $\mathbb{C}$ is the set of complex numbers $\{a+ib: a, b \in\mathbb{R}\}$ where $i = \sqrt{-1}$.
  \item $M_{mn}(\mathbb{R})$ is the set of $m\times n$ matrices with real coefficients.
  \item $\{:),\, :\!(\}$
\end{itemize}

We also have \underline{intervals} (subsets of $\mathbb{R}$)

\end{document}
